\documentclass[10pt,a4paper]{amsart}
\usepackage{amsmath,amssymb,amsthm,graphicx}
% Retained preamble for the Stacks Project chapter snapshot in this directory.
% The chapter's own preamble.tex is a copy of the shared Stacks Project preamble; the block
% below is sliced MECHANICALLY, line for line, from it. Nothing was added, removed or
% redefined; the wrapper class, the xr-hyper/hyperref wiring and the externaldocument list
% are omitted because the assembler supplies the document class.
%
% lines taken from preamble.tex: 3, 5, 9, 11, 12, 156, 157, 158, ... (51 lines total)
%
% package coverage restored from the chapter's own preamble (2026-10-13): xy, xy-2cell

\usepackage[all]{xy}
\xyoption{2cell}
\UseAllTwocells

\newenvironment{reference}{\comment}{\endcomment}
\newenvironment{slogan}{\comment}{\endcomment}
\newenvironment{history}{\comment}{\endcomment}
\theoremstyle{plain}
\newtheorem{theorem}[subsection]{Theorem}
\newtheorem{proposition}[subsection]{Proposition}
\newtheorem{lemma}[subsection]{Lemma}
\theoremstyle{definition}
\newtheorem{definition}[subsection]{Definition}
\newtheorem{example}[subsection]{Example}
\newtheorem{exercise}[subsection]{Exercise}
\newtheorem{situation}[subsection]{Situation}
\theoremstyle{remark}
\newtheorem{remark}[subsection]{Remark}
\newtheorem{remarks}[subsection]{Remarks}
\numberwithin{equation}{subsection}
\def\lim{\mathop{\mathrm{lim}}\nolimits}
\def\colim{\mathop{\mathrm{colim}}\nolimits}
\def\Spec{\mathop{\mathrm{Spec}}}
\def\Hom{\mathop{\mathrm{Hom}}\nolimits}
\def\Ext{\mathop{\mathrm{Ext}}\nolimits}
\def\SheafHom{\mathop{\mathcal{H}\!\mathit{om}}\nolimits}
\def\SheafExt{\mathop{\mathcal{E}\!\mathit{xt}}\nolimits}
\def\Sch{\mathit{Sch}}
\def\Mor{\mathop{\mathrm{Mor}}\nolimits}
\def\Ob{\mathop{\mathrm{Ob}}\nolimits}
\def\Sh{\mathop{\mathit{Sh}}\nolimits}
\def\NL{\mathop{N\!L}\nolimits}
\def\CH{\mathop{\mathrm{CH}}\nolimits}
\def\proetale{{pro\text{-}\acute{e}tale}}
\def\etale{{\acute{e}tale}}
\def\QCoh{\mathit{QCoh}}
\def\Ker{\mathop{\mathrm{Ker}}}
\def\Im{\mathop{\mathrm{Im}}}
\def\Coker{\mathop{\mathrm{Coker}}}
\def\Coim{\mathop{\mathrm{Coim}}}
\DeclareMathSymbol{\boxtimes}{\mathbin}{AMSa}{"02}
\def\QCohstack{\mathcal{QC}\!\mathit{oh}}
\def\Cohstack{\mathcal{C}\!\mathit{oh}}
\def\Spacesstack{\mathcal{S}\!\mathit{paces}}
\def\Quotfunctor{\mathrm{Quot}}
\def\Hilbfunctor{\mathrm{Hilb}}
\def\Curvesstack{\mathcal{C}\!\mathit{urves}}
\def\Polarizedstack{\mathcal{P}\!\mathit{olarized}}
\def\Complexesstack{\mathcal{C}\!\mathit{omplexes}}
\def\Pic{\mathop{\mathrm{Pic}}\nolimits}
\def\Picardstack{\mathcal{P}\!\mathit{ic}}
\def\Picardfunctor{\mathrm{Pic}}
\def\Deformationcategory{\mathcal{D}\!\mathit{ef}}
\usepackage{cleveref}
\crefname{originallemma}{original lemma}{original lemmas}
\Crefname{originallemma}{Original lemma}{Original lemmas}
\makeatletter
\newlabel{algebraic-definition-relative-representable-property}{{03YK}{source}{}{sourcepointer-1}{}}
\newlabel{algebraic-definition-relative-representable-property@cref}{{[originallemma][0][]03YK}{[1][0][]source}{}{}{}}
\makeatother
\crefname{originallemma}{original lemma}{original lemmas}
\Crefname{originallemma}{Original lemma}{Original lemmas}
\makeatletter
\newlabel{criteria-lemma-check-property-limit-preserving}{{07WI}{source}{}{sourcepointer-2}{}}
\newlabel{criteria-lemma-check-property-limit-preserving@cref}{{[originallemma][0][]07WI}{[1][0][]source}{}{}{}}
\makeatother
\crefname{originallemma}{original lemma}{original lemmas}
\Crefname{originallemma}{Original lemma}{Original lemmas}
\makeatletter
\newlabel{equation-smooth}{{07XS}{source}{}{sourcepointer-3}{}}
\newlabel{equation-smooth@cref}{{[originallemma][0][]07XS}{[1][0][]source}{}{}{}}
\makeatother
\crefname{originallemma}{original lemma}{original lemmas}
\Crefname{originallemma}{Original lemma}{Original lemmas}
\makeatletter
\newlabel{lemma-base-change-smooth}{{07XV}{source}{}{sourcepointer-4}{}}
\newlabel{lemma-base-change-smooth@cref}{{[originallemma][0][]07XV}{[1][0][]source}{}{}{}}
\makeatother
\crefname{originallemma}{original lemma}{original lemmas}
\Crefname{originallemma}{Original lemma}{Original lemmas}
\makeatletter
\newlabel{lemma-composition-smooth}{{07XU}{source}{}{sourcepointer-5}{}}
\newlabel{lemma-composition-smooth@cref}{{[originallemma][0][]07XU}{[1][0][]source}{}{}{}}
\makeatother
\crefname{originallemma}{original lemma}{original lemmas}
\Crefname{originallemma}{Original lemma}{Original lemmas}
\makeatletter
\newlabel{lemma-fibre-product-limit-preserving}{{07XM}{source}{}{sourcepointer-6}{}}
\newlabel{lemma-fibre-product-limit-preserving@cref}{{[originallemma][0][]07XM}{[1][0][]source}{}{}{}}
\makeatother
\crefname{originallemma}{original lemma}{original lemmas}
\Crefname{originallemma}{Original lemma}{Original lemmas}
\makeatletter
\newlabel{lemma-limit-preserving-algebraic-space}{{07XN}{source}{}{sourcepointer-7}{}}
\newlabel{lemma-limit-preserving-algebraic-space@cref}{{[originallemma][0][]07XN}{[1][0][]source}{}{}{}}
\makeatother
\crefname{originallemma}{original lemma}{original lemmas}
\Crefname{originallemma}{Original lemma}{Original lemmas}
\makeatletter
\newlabel{lemma-smooth-smooth}{{07XW}{source}{}{sourcepointer-8}{}}
\newlabel{lemma-smooth-smooth@cref}{{[originallemma][0][]07XW}{[1][0][]source}{}{}{}}
\makeatother
\crefname{originallemma}{original lemma}{original lemmas}
\Crefname{originallemma}{Original lemma}{Original lemmas}
\makeatletter
\newlabel{more-morphisms-lemma-lifting-along-artinian-at-point}{{02HX}{source}{}{sourcepointer-9}{}}
\newlabel{more-morphisms-lemma-lifting-along-artinian-at-point@cref}{{[originallemma][0][]02HX}{[1][0][]source}{}{}{}}
\makeatother
\crefname{originallemma}{original lemma}{original lemmas}
\Crefname{originallemma}{Original lemma}{Original lemmas}
\makeatletter
\newlabel{morphisms-lemma-enough-finite-type-points}{{02J4}{source}{}{sourcepointer-10}{}}
\newlabel{morphisms-lemma-enough-finite-type-points@cref}{{[originallemma][0][]02J4}{[1][0][]source}{}{}{}}
\makeatother
\title{The lifting property (equation-smooth) implies smoothness}
\author{The Stacks Project Authors}
\date{stacks-project:artin.tex:local-snapshot-20261012; GNU FDL 1.2 or later}
\begin{document}
\maketitle
\noindent\textit{Editorial note.} Statement span is lines 2344-2356 of sources/stacks-artin/pinned-original-source.tex; proof span is lines 2358-2402, the complete author proof (43 body lines).\par\medskip
\noindent\textit{Editorial note.} Navigation stubs: 10 cited-but-not-included labels are supplied as non-load-bearing pointers with their pinned tags where the snapshot tag index has them.\par\medskip
\noindent\textit{Editorial note.} extra\_packages is empty on purpose: the retained preamble already loads xy with options and a second declaration clashes in the assembler wrapper.\par\medskip
\noindent\textit{Editorial note.} Text is retained verbatim, including the author's own typos and placeholders.\par\medskip
\noindent\textit{Editorial note.} Queue index 20 of queue\_artin.json (pb=2113, stmt 2344-2356, proof 2358-2402).\par\medskip
\noindent\textit{Editorial note.} stacks-artin wave ar01 (2026-10-14): fresh adjudication accept per lane-s3/artin\_judgments.md; pool per Chief 10:5x direct dispatch (idle-first over S12's 01:5x note, SMC precedent); artin.tex 99 rows/67 with proof judged 46 accept / 21 reject / 5 AGP-overlap avoidance rows; patch-sd-envs contract spans scanned.\par\medskip
\noindent\textit{Allowed standard prerequisites (editorial).}\begin{itemize}
\item Explicit prerequisite boundary. The cited inputs are prerequisites; the argument in the retained span is not replaced by any cited result.
\end{itemize}\medskip
\ifdefined\hypertarget\hypertarget{sourcepointer-1}{}\fi
\noindent\textit{Source reference.} Original lemma 03YK: sources/stacks-artin/pinned-original-source.tex, source lines 2402--2402. Cited by the retained statement or proof but not part of the collected spans; supplied only as a navigation pointer. The argument does not depend on its content.\par\medskip
\ifdefined\hypertarget\hypertarget{sourcepointer-2}{}\fi
\noindent\textit{Source reference.} Original lemma 07WI: sources/stacks-artin/pinned-original-source.tex, source lines 2402--2402. Cited by the retained statement or proof but not part of the collected spans; supplied only as a navigation pointer. The argument does not depend on its content.\par\medskip
\ifdefined\hypertarget\hypertarget{sourcepointer-3}{}\fi
\noindent\textit{Source reference.} Original lemma 07XS: sources/stacks-artin/pinned-original-source.tex, source lines 2268--2275. Cited by the retained statement or proof but not part of the collected spans; supplied only as a navigation pointer. The argument does not depend on its content.\par\medskip
\ifdefined\hypertarget\hypertarget{sourcepointer-4}{}\fi
\noindent\textit{Source reference.} Original lemma 07XV: sources/stacks-artin/pinned-original-source.tex, source lines 2315--2321. Cited by the retained statement or proof but not part of the collected spans; supplied only as a navigation pointer. The argument does not depend on its content.\par\medskip
\ifdefined\hypertarget\hypertarget{sourcepointer-5}{}\fi
\noindent\textit{Source reference.} Original lemma 07XU: sources/stacks-artin/pinned-original-source.tex, source lines 2302--2307. Cited by the retained statement or proof but not part of the collected spans; supplied only as a navigation pointer. The argument does not depend on its content.\par\medskip
\ifdefined\hypertarget\hypertarget{sourcepointer-6}{}\fi
\noindent\textit{Source reference.} Original lemma 07XM: sources/stacks-artin/pinned-original-source.tex, source lines 1693--1706. Cited by the retained statement or proof but not part of the collected spans; supplied only as a navigation pointer. The argument does not depend on its content.\par\medskip
\ifdefined\hypertarget\hypertarget{sourcepointer-7}{}\fi
\noindent\textit{Source reference.} Original lemma 07XN: sources/stacks-artin/pinned-original-source.tex, source lines 1733--1742. Cited by the retained statement or proof but not part of the collected spans; supplied only as a navigation pointer. The argument does not depend on its content.\par\medskip
\ifdefined\hypertarget\hypertarget{sourcepointer-8}{}\fi
\noindent\textit{Source reference.} Original lemma 07XW: sources/stacks-artin/pinned-original-source.tex, source lines 2332--2338. Cited by the retained statement or proof but not part of the collected spans; supplied only as a navigation pointer. The argument does not depend on its content.\par\medskip
\ifdefined\hypertarget\hypertarget{sourcepointer-9}{}\fi
\noindent\textit{Source reference.} Original lemma 02HX: sources/stacks-artin/pinned-original-source.tex, source lines 2402--2402. Cited by the retained statement or proof but not part of the collected spans; supplied only as a navigation pointer. The argument does not depend on its content.\par\medskip
\ifdefined\hypertarget\hypertarget{sourcepointer-10}{}\fi
\noindent\textit{Source reference.} Original lemma 02J4: sources/stacks-artin/pinned-original-source.tex, source lines 2402--2402. Cited by the retained statement or proof but not part of the collected spans; supplied only as a navigation pointer. The argument does not depend on its content.\par\medskip
\begin{lemma}
\label{lemma-implies-smooth}
Let $S$ be a locally Noetherian scheme. Let $f : \mathcal{X} \to \mathcal{Y}$
be a $1$-morphism of categories fibred in groupoids over $(\Sch/S)_{fppf}$.
Assume
\begin{enumerate}
\item $f$ is representable by algebraic spaces,
\item $f$ satisfies (\ref{equation-smooth}),
\item $\mathcal{X} \to (\Sch/S)_{fppf}$ is limit preserving on objects, and
\item $\mathcal{Y}$ is limit preserving.
\end{enumerate}
Then $f$ is smooth.
\end{lemma}

\begin{proof}
The key ingredient of the proof is More on Morphisms, Lemma
\ref{more-morphisms-lemma-lifting-along-artinian-at-point}
which (almost) says that a morphism of schemes of finite type over $S$
satisfying (\ref{equation-smooth}) is a smooth morphism. The other
arguments of the proof are essentially bookkeeping.

\medskip\noindent
Let $V$ be a scheme over $S$ and let $y$ be an object of $\mathcal{Y}$ over
$V$. Let $Z$ be an algebraic space representing the $2$-fibre product
$\mathcal{Z} = \mathcal{X} \times_{f, \mathcal{X}, y} (\Sch/V)_{fppf}$.
We have to show that the projection morphism $Z \to V$ is smooth, see
Algebraic Stacks, Definition
\ref{algebraic-definition-relative-representable-property}.
In fact, it suffices to do this when $V$ is an affine scheme
locally of finite presentation over $S$, see
Criteria for Representability, Lemma
\ref{criteria-lemma-check-property-limit-preserving}.
Then $(\Sch/V)_{fppf}$ is limit preserving by
Lemma \ref{lemma-limit-preserving-algebraic-space}.
Hence $Z \to S$ is locally of finite presentation by
Lemmas \ref{lemma-fibre-product-limit-preserving} and
\ref{lemma-limit-preserving-algebraic-space}.
Choose a scheme $W$ and a surjective \'etale morphism $W \to Z$.
Then $W$ is locally of finite presentation over $S$.

\medskip\noindent
Since $f$ satisfies (\ref{equation-smooth}) we see that so does
$\mathcal{Z} \to (\Sch/V)_{fppf}$, see Lemma \ref{lemma-base-change-smooth}.
Next, we see that $(\Sch/W)_{fppf} \to \mathcal{Z}$ satisfies
(\ref{equation-smooth}) by Lemma \ref{lemma-smooth-smooth}.
Thus the composition
$$
(\Sch/W)_{fppf} \to \mathcal{Z} \to (\Sch/V)_{fppf}
$$
satisfies (\ref{equation-smooth}) by Lemma \ref{lemma-composition-smooth}.
More on Morphisms, Lemma
\ref{more-morphisms-lemma-lifting-along-artinian-at-point}
shows that the composition $W \to Z \to V$ is smooth at every finite type
point $w_0$ of $W$. Since the smooth locus is open we conclude
that $W \to V$ is a smooth morphism of schemes by
Morphisms, Lemma \ref{morphisms-lemma-enough-finite-type-points}.
Thus we conclude that $Z \to V$ is a smooth morphism
of algebraic spaces by definition.
\end{proof}

\end{document}
